% ---
% title: Calls to a help desk
% tags: exponential, distribution-function
% ---

\begin{problem}
	The length $X$, in minutes, of a call to a bank's help desk is
	exponential with rate $\lambda=1/10$, so its density is
	\[
		f(x)=\frac1{10}\,e^{-x/10}\quad\text{for } x\ge0,
	\]
	and $f(x)=0$ for $x<0$.
	(a) Find the probability that a call lasts more than 10 minutes.
	(b) Find the probability that a call lasts between 10 and 20 minutes.
	(c) Find the probability that a call is shorter than 5 minutes.
\end{problem}

\begin{solution}
	\idea First find the distribution function $F(x)=P(X\le x)$ by
	integrating the density once. Then every probability in the exercise is a
	value of $F$ or a difference of two values of $F$; no new integral is
	needed.

	\step{The distribution function.} For $x<0$ there is no density to the
	left of $x$, so $F(x)=0$. For $x\ge0$ we integrate the density from $0$
	to $x$. An antiderivative of $\tfrac1{10}e^{-t/10}$ is $-e^{-t/10}$
	(differentiate it to check: the chain rule brings down the factor
	$-\tfrac1{10}$). So
	\[
		F(x)=\int_0^x \frac1{10}\,e^{-t/10}\,dt
		=\Big[{-e^{-t/10}}\Big]_0^x
		=-e^{-x/10}-\big(-e^{0}\big)
		=1-e^{-x/10}.
	\]
	The complement is even simpler: the probability that a call lasts
	more than $x$ minutes is
	\[
		P(X>x)=1-F(x)=e^{-x/10}\qquad(x\ge0).
	\]
	The picture shows the density; the shaded area is the probability in (b).
	\begin{center}
		\begin{tikzpicture}[x=0.22cm,y=28cm,>=stealth]
			\fill[teal!35] (10,0) -- plot[domain=10:20,samples=40]
			(\x,{0.1*exp(-\x/10)}) -- (20,0) -- cycle;
			\draw[->] (-1,0)--(47,0) node[right] {$x$};
			\draw[->] (0,-0.005)--(0,0.125) node[above] {$f(x)$};
			\draw[navy,very thick] plot[domain=0:45,samples=60]
			(\x,{0.1*exp(-\x/10)});
			\foreach \x in {5,10,20,30,40}
			\draw (\x,0.002)--(\x,-0.002) node[below] {\footnotesize$\x$};
			\draw (0.4,0.1)--(-0.4,0.1) node[left] {\footnotesize$\frac1{10}$};
			\node[below left] at (0,0) {\footnotesize$0$};
		\end{tikzpicture}
	\end{center}

	\step{(a)} ``More than 10 minutes'' is the event $X>10$. By the formula
	above,
	\[
		P(X>10)=e^{-10/10}=e^{-1}\approx0.368.
	\]
	So a bit more than one call in three lasts longer than 10 minutes.

	\step{(b)} The event $X\le20$ splits into the two disjoint events
	$X\le10$ and $10<X\le20$. So $F(20)=F(10)+P(10<X\le20)$, that is,
	\[
		P(10<X\le20)=F(20)-F(10)
		=\big(1-e^{-2}\big)-\big(1-e^{-1}\big)
		=e^{-1}-e^{-2}\approx0.3679-0.1353=0.2325.
	\]
	$X$ is continuous, so a single point has probability $0$: it does not
	matter whether we write $10<X<20$ or $10<X\le20$. Hence
	$P(10<X<20)=e^{-1}-e^{-2}\approx0.233$.
	This is the shaded area in the picture.

	\step{(c)} ``Shorter than 5 minutes'' is the event $X<5$. Again
	$P(X=5)=0$, so $P(X<5)=P(X\le5)=F(5)$:
	\[
		P(X<5)=1-e^{-5/10}=1-e^{-0.5}\approx1-0.6065=0.3935\approx0.393.
	\]

	\begin{answer}
		(a) $e^{-1}\approx0.368$ \quad
		(b) $e^{-1}-e^{-2}\approx0.233$ \quad
		(c) $1-e^{-0.5}\approx0.393$
	\end{answer}
	\begin{remark}
		\textbf{Common mistake:} using $10$ instead of $1/10$ as the
		rate. The rate $\lambda$ appears in the exponent as $e^{-\lambda x}$;
		here $\lambda=1/10$ per minute, so the exponent is $-x/10$. With
		$e^{-10x}$ the probability in (a) would be $e^{-100}$, practically
		$0$, which makes no sense for calls that often last several minutes.
		A useful check: for an exponential distribution with rate $\lambda$,
		the probability of lasting more than $1/\lambda$ (here 10 minutes) is
		always $e^{-1}\approx0.368$.
	\end{remark}
\end{solution}
